\documentclass[11pt,reqno]{amsart}
\usepackage[T1]{fontenc}
\usepackage{lmodern,amsmath,amssymb,amsthm,mathtools,microtype}
\usepackage[margin=1.1in]{geometry}
\usepackage[hidelinks]{hyperref}
\allowdisplaybreaks[2]
\setlength{\emergencystretch}{2em}
\numberwithin{equation}{section}
\newtheorem{theorem}{Theorem}[section]
\newtheorem{lemma}[theorem]{Lemma}
\newtheorem{proposition}[theorem]{Proposition}
\newtheorem{corollary}[theorem]{Corollary}
\theoremstyle{remark}\newtheorem{remark}[theorem]{Remark}
\newcommand{\R}{\mathbb R}\newcommand{\Z}{\mathbb Z}\newcommand{\E}{\mathbb E}
\newcommand{\dd}{\,\mathrm d}\newcommand{\one}{\mathbf1}\newcommand{\T}{\mathbb T}
\newcommand{\nuR}{\nu}\newcommand{\av}[1]{\int #1\dd\nu}
\title[Two-collision counting response]{Two-collision counting response in a periodic Lorentz gas}
\author{Qian Qi}
\date{September 8, 2026}
\subjclass[2020]{37D50, 37A50, 60G55}
\keywords{Lorentz gas, collision counts, moving boundaries, suspension flow, response}
\begin{document}
\begin{abstract}
For the periodic triangular Lorentz gas with disk radius in $[9/20,47/100]$, we determine the full collision-count law over the physical time interval of length $11/100$. The window admits two collisions with positive probability for every radius, rather than excluding the second collision. Stationary flux identifies the two-collision probability with a positive-part integral of the genuine free-flight roof. A geometric short-flight estimate makes the entire support of this integral uniformly transverse, and its terminal level is a regular level throughout the radius interval. We prove smooth radius response of all three count probabilities, strict response of the two-collision probability, and explicit formulae for the moving-level terms at every finite order. No independence of consecutive collisions is imposed. We also give the exact collision-tail formula for arbitrary bounded windows. The results concern count statistics under the full equilibrium preparation; they do not assert a long-time susceptibility theorem or response of the entire two-collision path distribution.
\end{abstract}
\maketitle

\section{The mechanical family and the result}
Let $e_1=(1,0)$, $e_2=(1/2,\sqrt3/2)$, and $\Lambda=\Z e_1+\Z e_2$. On $\T^2_\Lambda=\R^2/\Lambda$ remove the disk of radius $R$, where
\begin{equation}\label{eq:interval}
 I=[9/20,47/100],\qquad A_R=\sqrt3/2-\pi R^2,\qquad T=11/100.
\end{equation}
The particle has unit speed and undergoes Euclidean elastic reflection. Its preparation is the entire normalized phase volume
\[
 \mu_R(\dd q,\dd\vartheta)=\frac{\one_{Q_R}(q)}{2\pi A_R}\dd q\dd\vartheta,
 \qquad Q_R=\T^2_\Lambda\setminus\overline B(0,R).
\]
Let $N_T$ count collisions in $(0,T]$. Endpoint and tangent conventions alter no probabilities below.

On the outgoing collision section write $q=R n_\alpha$ and
\[
 v=\cos\phi\,n_\alpha+\sin\phi\,t_\alpha,
 \quad n_\alpha=(\cos\alpha,\sin\alpha),\quad
 t_\alpha=(-\sin\alpha,\cos\alpha),\quad |\phi|<\pi/2.
\]
The normalized section measure and its mean physical roof are
\begin{equation}\label{eq:flux}
 \dd\nu=\frac{\cos\phi}{4\pi}\dd\alpha\dd\phi,
 \qquad \overline\tau_R=\int\tau_R\dd\nu=\frac{A_R}{2R},
 \qquad \lambda_R=\overline\tau_R^{-1}.
\end{equation}
Here $\tau_R$ is the actual next collision time, not a prescribed symbolic clock. Put
\begin{equation}\label{eq:A}
 J(R)=\int(T-\tau_R)_+\dd\nu.
\end{equation}

\begin{theorem}\label{thm:main}
For the family \eqref{eq:interval}, $N_T\in\{0,1,2\}$ almost surely and
\begin{equation}\label{eq:law}
 p_2=\lambda_R J(R),\qquad
 p_1=\lambda_R T-2\lambda_R J(R),\qquad
 p_0=1-\lambda_R T+\lambda_R J(R).
\end{equation}
All three probabilities are strictly positive and extend smoothly to a neighborhood of $I$. Every fixed finite derivative is bounded uniformly on $I$. Moreover
\[
 p_2(R)\ge 10^{-9},\qquad p_2'(R)>0\quad(R\in I).
\]
For each nearest lattice center $e$, let $\tau_{e,R}$ be the smooth entry root on the short-flight branch, and let $D_e(R)=\{\tau_{e,R}<T\}$ in section coordinates. Then
\begin{align}\label{eq:firstsecond}
 J'(R)&=-\sum_e\int_{D_e(R)}\partial_R\tau_{e,R}\dd\nu,\\
 J''(R)&=-\sum_e\int_{D_e(R)}\partial_R^2\tau_{e,R}\dd\nu
 +\sum_e\int_{\tau_{e,R}=T}
 \frac{(\partial_R\tau_{e,R})^2\cos\phi}{4\pi|\nabla_{\alpha,\phi}\tau_{e,R}|}\dd\mathcal H^1.
\end{align}
The formulas include the preparation normalizer through \eqref{eq:law}. They do not require a grazing cutoff on the initial equilibrium preparation.
\end{theorem}

The integral \eqref{eq:A} is a deterministic flux integral. We do not introduce a renewal assumption. The stationary counting identities used in the proof are elementary suspension identities, closely related to the Palm description of entry times \cite{Marklof}. The argument combines this stationary identity with regularity of the complete short-flight region. In particular, the one-event computation is not simply iterated.

\section{Flux and exact collision tails}
\subsection{Geometry and equilibrium suspension}
The shortest nonzero lattice vector has length one. Distinct disks are separated by at least
\begin{equation}\label{eq:gap}
 d_*=1-2\max I=3/50.
\end{equation}
An outgoing ray cannot hit the same convex disk again. Thus successive collisions have gap at least $d_*$. There can be no finite-time collision accumulation. Grazing directions from a fixed exterior position form a finite or countable null set; regular collision maps preserve flux, and their finite preimages are null. Removing their countable union gives the full-measure regular dynamics used here.

For completeness the family has finite horizon. For a primitive lattice direction $w$, parallel lattice rows have spacing $(\sqrt3/2)/|w|\le\sqrt3/2$. A parallel line has distance at most $\sqrt3/4<9/20$ from a row, and eventually passes through a disk. An irrational-direction line is dense on the torus. If arbitrarily long collision-free segments existed at the smallest radius, a convergent sequence of starting positions and directions would yield an infinite ray avoiding every open disk, a contradiction. Compactness gives a common finite roof upper bound. Enlarging the disks preserves this upper bound.

A boundary flow tube has volume $\cos\phi\dd s\dd\phi\dd t$, where $\dd s=R\dd\alpha$. Its total section flux is $4\pi R$, while total unit-speed phase volume is $2\pi A_R$. Consequently
\[
 2\pi A_R=4\pi R\int\tau_R\dd\nu.
\]
Reflection reverses normal velocity and preserves flux. Free flight preserves phase volume. The collision map $B_R$ therefore preserves $\nu$, and \eqref{eq:flux} follows. The flow preparation is the suspension measure
\begin{equation}\label{eq:suspension}
 \frac{\dd\nu(x)\dd s}{\overline\tau_R},\qquad 0\le s<\tau_R(x).
\end{equation}
No mixing property enters this statement.

\subsection{A general counting identity}
Let $B$ be an invertible probability-preserving map with roof $\tau\ge d_*>0$ and $\overline\tau=\int\tau\dd\nu<\infty$. Write $S_j\tau(x)=\sum_{h=0}^{j-1}\tau(B^hx)$, with $S_0=0$.
\begin{proposition}\label{prop:tails}
Under the stationary suspension preparation, for every $k\ge1$ and $t\ge0$,
\begin{equation}\label{eq:tails}
 \Pr(N_t\ge k)=\frac1{\overline\tau}\int
 \bigl[(t-S_{k-1}\tau)_+-(t-S_k\tau)_+\bigr]\dd\nu.
\end{equation}
In particular $\E N_t=t/\overline\tau$. These formulas hold without independent roofs or a Markov hypothesis.
\end{proposition}
\begin{proof}
From a suspension point $(x,s)$, the $k$th future collision occurs after
$S_k\tau(x)-s$. The length of the subset of $[0,\tau(x))$ for which this is at most $t$ is
\[
 (t-S_{k-1}\tau(Bx))_+-(t-S_k\tau(x))_+.
\]
This follows by intersecting the interval $s\ge S_k\tau(x)-t$ with $[0,\tau(x))$. Integrate using \eqref{eq:suspension}. Invariance of $\nu$ replaces $Bx$ by $x$ in the first term, proving \eqref{eq:tails}. Because $S_j\tau\ge jd_*$, only finitely many tail terms are nonzero for fixed $t$. Summing in $k$ telescopes to $t/\overline\tau$.
\end{proof}

When $T<2d_*$, the second tail is $\lambda_R J(R)$ and every third tail is zero. Also the first tail is $\lambda_R(T-J(R))$. Taking differences proves \eqref{eq:law}. Notice the correction to the former one-event identity:
\begin{equation}\label{eq:correction}
 \Pr(N_T\ge1)=\lambda_RT-p_2,
 \qquad \E N_T=\lambda_RT.
\end{equation}
Once a second event is possible, event probability and expected count are different.

More generally, with $A_j(t)=\int(t-S_j\tau)_+\dd\nu$, the probability of exactly $k\ge1$ collisions is
\[
 \overline\tau^{-1}(A_{k-1}-2A_k+A_{k+1}).
\]
This is an exact decomposition at every bounded horizon. It does not supply the parameter regularity of the roof iterates $S_j\tau$.

\section{The complete short-flight region is regular}
Let $q=Rn_\alpha$, $v=\cos\phi\,n_\alpha+\sin\phi\,t_\alpha$ and suppose the next impact is at $q'=e+Rn_1=q+\tau v$. Set $c_0=n_\alpha\cdot v>0$ and $c_1=-n_1\cdot v>0$.
\begin{lemma}\label{lem:transverse}
Every next flight with $\tau\le T$ ends at one of the six nearest lattice centers. Both endpoint cosines satisfy
\begin{equation}\label{eq:cos}
 c_0,c_1\ge\frac{1-2R-T^2}{2RT}\ge\frac{479}{1034}>\frac{46}{100}.
\end{equation}
All entry roots with length at most $T$ are therefore in finitely many smooth branch neighborhoods. No other disk obstructs such a root.
\end{lemma}
\begin{proof}
The triangle inequality gives $|e|\le2R+T\le1.05<\sqrt3$. The next possible lattice length after one is $\sqrt3$, so $|e|=1$. Since $|q'|\ge1-R$,
\[
 (1-R)^2\le |q'|^2=R^2+2R\tau c_0+\tau^2.
\]
The analogous identity $|q-e|^2=R^2+2R\tau c_1+\tau^2$ proves the same bound for $c_1$. The function $(1-2R-\tau^2)/(2R\tau)$ decreases in both variables in these ranges. Its value at $R=47/100$, $\tau=11/100$ is $479/1034$. A hypothetical earlier different obstacle followed by the candidate obstacle along the straight segment would require at least two inter-obstacle gaps, exceeding $T$; a repeated encounter with the initial disk is excluded by convexity. Thus every outgoing entry root of length at most $T$ is the first actual hit. The implicit root denominator is $-2Rc_1$, uniformly separated from zero. Compactness and the six labels give the claimed neighborhoods.
\end{proof}

This lemma does not say that all first collisions from arbitrary phase points are transverse. In the stationary count identity, the integral concerns the complete intercollision flight following a section point. Grazing first events remain in the equilibrium ensemble. The count statistic has eliminated their individual decoding through an exact identity, not by assigning them zero probability.

On a branch let $d=e-Rn_\alpha$, $L=d\cdot v$, and $b=d\cdot v^\perp$. Then
\begin{equation}\label{eq:root}
 \tau_{e,R}=L-\sqrt{R^2-b^2},\qquad
 \partial_R\tau_{e,R}=\frac{1-n_1\cdot n_\alpha}{n_1\cdot v}<0.
\end{equation}
The root formula solves $|q+tv-e|^2=R^2$ on its entry branch. Differentiating that equation at fixed $(\alpha,\phi)$ gives the derivative formula; $v$ is independent of $R$ in these coordinates. If its numerator vanished, $n_1=n_\alpha$, contrary to $n_\alpha\cdot v>0>n_1\cdot v$. It is thus strictly positive.

\begin{lemma}\label{lem:regularlevel}
The level $\tau_{e,R}=T$ is regular in the two section variables. The regularity is uniform for $R\in I$ over its complete level set.
\end{lemma}
\begin{proof}
Use $(\alpha,\vartheta)$, where $\vartheta=\alpha+\phi$ is the direction angle. This is a fixed invertible linear change of angular coordinates. Implicit differentiation gives
\[
 \partial_\vartheta\tau=-\tau\frac{n_1\cdot v^\perp}{n_1\cdot v},
 \qquad \partial_\alpha\tau\big|_\vartheta=-R\frac{n_1\cdot t_\alpha}{n_1\cdot v}.
\]
If both vanished, $n_1$ would be parallel to both $v$ and $n_\alpha$. The incoming and outgoing signs force $v=n_\alpha$ and $n_1=-n_\alpha$. The chord equation becomes $e=(2R+\tau)n_\alpha$, so $\tau=1-2R\le1/10<T$. This contradicts the level equation. The short-flight support is compactly separated from tangency by Lemma~\ref{lem:transverse}; the level sets form a compact subset of the finitely many branch neighborhoods. Continuity gives a positive lower gradient bound there. All inequalities are strict on $I$, so the result extends to an open radius neighborhood.
\end{proof}

\section{Moving-level response at every finite order}
\subsection{Smoothness of the positive-part integral}
The complete active set $\{\tau_R\le T\}$ is contained in the regular branch neighborhoods of Lemma~\ref{lem:transverse}. Choose a smooth partition of unity on a neighborhood of its compact union in parameter--section space. Terms away from $\tau=T$ are ordinary smooth parameter integrals: the positive part either vanishes or is $T-\tau$. Near the level, the parametric implicit function theorem supplies local coordinates $(R,y,u)$ with $u=\tau_{e,R}$, by Lemma~\ref{lem:regularlevel}. The corresponding integral is
\begin{equation}\label{eq:flatten}
 \int (T-u)_+a(R,y,u)\dd y\dd u,
\end{equation}
where $a$ is smooth and compactly supported in a common chart. Differentiation in $R$ acts only on $a$. Every fixed finite derivative is integrably bounded uniformly on a compact neighborhood of $I$. A finite covering proves $J\in C^\infty$ with the asserted finite-order bounds. The argument includes derivatives of the partition functions; in their sum they cancel by the partition identity. No analytic or order-uniform factorial bound is asserted.

\subsection{Explicit derivatives}
The first derivative has no level term because the factor $T-\tau$ vanishes on the moving level. Differentiating it once more gives \eqref{eq:firstsecond} by the coarea formula. More explicitly,
\[
 \partial_R\one_{\{\tau_R<T\}}=-\tau_R'\delta(T-\tau_R),
\]
and
\[
 \int f\delta(T-\tau_R)\dd\nu
 =\int_{\tau_R=T}\frac{f\cos\phi}{4\pi|\nabla\tau_R|}\dd\mathcal H^1.
\]
Both identities can be verified directly in \eqref{eq:flatten}; their use does not assume a product of unrelated singular distributions.

Define the partial Bell polynomials by
\[
 B_{r,k}(x_1,\ldots)=\sum_{\substack{\sum m_j=k\\\sum jm_j=r}}
 \frac{r!}{\prod_jm_j!}\prod_j(x_j/j!)^{m_j}.
\]
For $r\ge1$ the full finite-order formula on the regular neighborhoods is
\begin{align}\label{eq:alljets}
 J^{(r)}(R)=\sum_e\int\Bigl[&-\one_{\{\tau_{e,R}<T\}}\tau_{e,R}^{(r)}\\
 &+\sum_{k=2}^r(-1)^k B_{r,k}(\tau_{e,R}',\ldots)
 \delta^{(k-2)}(T-\tau_{e,R})\Bigr]\dd\nu.\notag
\end{align}
The formula is understood with local cutoffs equal to one on the relevant support. It is independent of their choice. The superscript on $\delta$ differentiates its scalar argument. To prove \eqref{eq:alljets}, compose the distribution $(T-u)_+$ with $u=\tau_R$ in the flattened charts. Its first derivative in $u$ is $-\one_{u<T}$ and its $k$th derivative is $(-1)^k\delta^{(k-2)}(T-u)$ for $k\ge2$. The finite chain rule gives the formula. Alternatively mollify the scalar positive part, apply the ordinary chain rule, and pass to distributions using the regular-level coordinates. This proves every prescribed finite order.

The roof jets themselves have an explicit algebraic recursion. Put $h_R=\sqrt{R^2-b_R^2}>0$ on a short-flight branch. Then
\[
 h_R^{(j)}=\frac{(R^2-b_R^2)^{(j)}-
 \sum_{k=1}^{j-1}\binom jk h_R^{(k)}h_R^{(j-k)}}{2h_R},
 \qquad \tau_R^{(j)}=L_R^{(j)}-h_R^{(j)}.
\]
Here $L_R$ and $b_R$ are affine in $R$, and $h_R=Rc_1$ has a positive uniform lower bound by \eqref{eq:cos}. Coarea derivatives are computed from the smooth density in \eqref{eq:flatten}. Finally
\[
 \lambda_R^{(j)}=\frac{(2R)^{(j)}-
 \sum_{k=1}^{\min(j,2)}\binom jk A_R^{(k)}\lambda_R^{(j-k)}}{A_R}
\]
and the ordinary Leibniz rule in \eqref{eq:law} determine every count-probability jet. This specifies the operations and all denominators.

\section{Positivity, strict response, and statistical consequences}
\subsection{A radius-uniform two-collision witness}
Consider the outgoing section square $|\alpha|\le10^{-3}$, $|\vartheta|\le10^{-3}$ facing the disk at $e_1$. It has $|\phi|\le2\cdot10^{-3}$. For $d=e_1-Rn_\alpha$,
\[
 L=d\cdot v\le1-R+10^{-6},\qquad
 |b|\le1.94\cdot10^{-3}.
\]
These bounds follow from $|\sin x|\le|x|$ and $1-\cos x\le x^2/2$: for example
$L=\cos\vartheta-R\cos(\alpha-\vartheta)$, and
$|b|\le |\sin\vartheta|+R|\sin(\alpha-\vartheta)|\le1.94\cdot10^{-3}$. Consequently,
\[
 0\le R-\sqrt{R^2-b^2}
 =\frac{b^2}{R+\sqrt{R^2-b^2}}<5\cdot10^{-6}.
\]
Using $L\le1-R+10^{-6}$ yields
$\tau_{e_1,R}<1-2R+6\cdot10^{-6}<0.10001$.
The root is outgoing from the initial disk, incoming at the target, and physical by Lemma~\ref{lem:transverse}'s no-obstruction argument. Hence $T-\tau>9/1000$ throughout this square. Its $\nu$ mass exceeds $1/4{,}000{,}000$, since its area is $4\cdot10^{-6}$ and $\cos\phi\ge1-2\cdot10^{-6}$. Also $\lambda_R>9/10$ because $A_R<1$. Therefore
\[
 p_2=\lambda_RJ(R)>
 \frac9{10}\frac9{1000}\frac1{4{,}000{,}000}>10^{-9}.
\]
Since $\tau_R'<0$ throughout every active short-flight branch, \eqref{eq:firstsecond} and this positive-mass square give $J'(R)>0$. Furthermore
\[
 \lambda_R'=\frac{2(A_\Lambda+\pi R^2)}{A_R^2}>0.
\]
Thus $p_2'=\lambda_R'J+\lambda_RJ'>0$.

For $p_1$, use $J\le T-d_*$ to obtain
\[
 p_1\ge\lambda_R(2d_*-T)>0.
\]
For $p_0$, $p_0\ge1-\lambda_RT$. A coarse rational bound suffices: $A_\Lambda>43/50$, $\pi<22/7$, and $R\le47/100$ imply $A_R>0.16$, while $2RT<0.104$. Consequently $\lambda_RT<0.65<1$. This completes the proof of Theorem~\ref{thm:main}.

\subsection{Count generating functions and identification}
\begin{corollary}\label{cor:source}
The count partition function is
\[
 Z_R(q)=p_0+p_1e^q+p_2e^{2q}
 =1+\lambda_RT(e^q-1)+\lambda_RJ(R)(e^q-1)^2.
\]
It is smooth in $R$, analytic in real $q$, and positive. Its log-source Hessian is strictly positive. Each target mean in $(0,2)$ has a unique finite conjugate source. The binary observation $B=\one_{\{N_T=2\}}$ identifies $R$ and has Fisher information
\[
 \mathcal I_B(R)=\frac{[p_2'(R)]^2}{p_2(R)(1-p_2(R))}>0.
\]
Its square-root probability vector is differentiable in quadratic mean.
\end{corollary}
\begin{proof}
Insert \eqref{eq:law} and collect powers of $e^q-1$. All probabilities are positive, so the tilted count distribution has nonzero variance. Its mean tends to zero and two as $q$ tends to opposite infinities, and is strictly increasing. The implicit function theorem gives local parameter derivatives of the selected source. Strict monotonicity of $p_2$ gives identification. Differentiating the two Bernoulli probabilities yields the information formula. Smoothness and positivity on a compact radius neighborhood give bounded second derivatives of their square roots, proving the quadratic-mean assertion by Taylor's formula. These are finite-experiment consequences, not a claim of a new general statistical principle.
\end{proof}

\section{What the next mechanical estimate must address}
The second-collision probability is obtained on the full natural preparation and includes a nonzero terminal-level contribution to its second derivative. The method does not construct the full joint two-impact path-law response: earlier and later impact directions are correlated, and sampling velocities across event times needs its own decoder analysis. It also does not supply a derivative bound on $S_j\tau_R$ uniform in $j$ in \eqref{eq:tails}. Those are distinct next estimates.

For A2, a concrete continuation is a source-dependent two-impact record assembly with fixed physical sampling tests, followed by a uniform bound on the complete finite-itinerary response. Long-time pressure or covariance response must use a genuinely summable version of that bound. Existing functional-analytic perturbation theory for the Lorentz gas \cite{DemersZhang} provides relevant comparison, not an automatic all-order theorem for these moving tests. The present result is an incremental mechanical input, not a replacement for that long-time task.

\begin{thebibliography}{9}
\bibitem{Marklof} J. Marklof, \emph{Entry and return times for semi-flows}, arXiv:1605.02715 (2016).\par\url{https://arxiv.org/abs/1605.02715}
\bibitem{DemersZhang} M. F. Demers and H.-K. Zhang, \emph{A functional analytic approach to perturbations of the Lorentz gas}, Comm. Math. Phys. (2013); arXiv:1210.1261.\par\url{https://arxiv.org/abs/1210.1261}
\bibitem{DYN} Q. Qi, \emph{Response and fluctuations of moving-cut hyperbolic path ensembles}, manuscript of September 6, 2026, Appendix C. Source snapshot and SHA-256 are recorded in the accompanying audit. This note rederives its flux input and extends the count window; it does not import its long-time theorems into billiards.
\end{thebibliography}
\end{document}
